\medskip\noindent\textbf{T1 and T2: Exact proof transport (from the disclosed prefix).}
Let $\pi=(\pi_1,\ldots,\pi_k)$ be a valid finite $F$-proof of $\varphi$ from $\Gamma$.
Hypothesis H1 carries each hypothesis line and each legal inference to a valid corresponding
line in $F^t$. It follows that $t(\pi)$ is an $F^t$-proof of $t(\varphi)$ from $t[\Gamma]$
with exactly $k$ lines. The inverse recoding gives the reverse implication and has the
same length-preserving property. Consequently theoremhood is preserved and the attainable
proof-length sets agree, including when no proof exists. Their minima and all exact layers
therefore agree. This proves T1 and T2.

\medskip\noindent\textbf{T3: Certified eventual emission (from the disclosed prefix).}
Fix a provable target $\varphi$ and a finite valid proof code $c$. By H3, the exhaustive
code enumeration assigns $c$ a finite position $j$. The protected enumerator receives
sufficient recurring service to reach position $j$ at a finite stage, unless an earlier
successful halt has already established emission. H2 recognizes $c$ as valid, and the
specified emission mechanism records it. Thus the required finite certificate and emission
stage exist. No successful guidance heuristic is required. This proves T3 under H2--H3.

\medskip\noindent\textbf{T4: Conditional candidate-rank calculation (new continuation).}
Fix a benchmark task $i$ and abbreviate $g=g_i$ and $p=p_i$. Because the baseline
uses a uniformly random order on $N$ candidates, its single accepting candidate has
uniform rank, so
\[
\mathbb E[R_{\mathsf B,i}]=\frac{N+1}{2}.
\]
For guided search, condition on whether the accepting certificate belongs to the target-blind
neighborhood $G_i$ of size $g$. If it belongs, its conditional expected rank is $(g+1)/2$.
If it does not, the guided method first examines and rejects all $g$ neighborhood
candidates, then encounters the accepting candidate at an expected rank $(N-g+1)/2$
within the complement. Therefore its overall rank is
\begin{align*}
\mathbb E[R_{\mathsf H,i}]
 &=p\frac{g+1}{2}+(1-p)\left(g+\frac{N-g+1}{2}\right)\\
 &=\frac{N+g+1-pN}{2}.
\end{align*}
Subtracting this from the baseline expectation yields
\[
\mathbb E[R_{\mathsf B,i}-R_{\mathsf H,i}]
=\frac{p_iN-g_i}{2}\ge\frac{p_0N-g_0}{2},
\]
where the last inequality uses $p_i\ge p_0$ and $g_i\le g_0$.
By linearity of expectation, the gross rank saving on $q$ tasks is at least
$q(p_0N-g_0)/2$, whether or not the tasks are independent.
Under the common charged comparison model H6, identical common costs cancel only
if they were genuinely incurred equally. H5 charges the additional guided work and bounds it by
\[
D_q\le a_0NL_N+q[(a_1\beta_0\Delta_0+a_2g_0)L_N+a_3].
\]
Subtracting this upper bound from the expected gross benefit establishes
\begin{equation}\label{eq:maincandidate}
\mathbb E[C_{\mathsf B}-C_{\mathsf H}]
\ge\frac q2(p_0N-g_0)-a_0NL_N
-q[(a_1\beta_0\Delta_0+a_2g_0)L_N+a_3].
\end{equation}
When the right side is positive, the guided method has strictly smaller expected
fully charged cost. This is precisely T4.

\medskip\noindent\textbf{Amortized threshold (new continuation).}
Set $A=a_1\beta_0\Delta_0+a_2g_0$ and $L=L_N=\lceil\log_2(N+1)\rceil$.
The right side of \eqref{eq:maincandidate} becomes
\[
q\left[\frac{p_0N-g_0}{2}-AL-a_3\right]-a_0NL.
\]
Write $b_N=(p_0N-g_0)/2-AL-a_3$. Since $p_0>0$ and the other parameters are fixed,
$b_N=(p_0/2)N-O(\log N)$, so $b_N>0$ for all sufficiently large $N$.
For $q\ge KL$ with $K>2a_0/p_0$, we may therefore bound the expression below by
\begin{align*}
KLb_N-a_0NL
&=L\left[K\left(\frac{p_0N-g_0}{2}-AL-a_3\right)-a_0N\right]\\
&=L\left[\left(\frac{Kp_0}{2}-a_0\right)N
-KAL-\frac{Kg_0}{2}-Ka_3\right].
\end{align*}
The coefficient $(Kp_0/2-a_0)$ is strictly positive, while $L=O(\log N)$;
therefore the positive linear term eventually dominates the logarithmic and
constant terms. Choose $N_0$ large enough that the bracket and $b_N$
are positive for all $N\ge N_0$. The lower bound is then strictly positive
for every integer $q\ge KL_N$. This proves the asymptotic assertion. \qedhere
